\tarjetavf{Pronósticos - Errores}{I2}{%
Los errores sistemáticos asociados a pronósticos son aquellos que no pueden ser explicados con ningún modelo.
}{%
\textbf{FALSO.} Esos son los errores aleatorios (ruido blanco). Los errores sistemáticos representan un sesgo que sí podría explicarse ajustando el modelo.
}

\tarjetavf{Inventarios (EOQ)}{I1}{%
Entre los supuestos del modelo EOQ está que el costo de emitir una orden (Setup) disminuye al aumentar la cantidad ordenada.
}{%
\textbf{FALSO.} El costo de emitir una orden ($S$) es constante e independiente del tamaño del lote.
}

\tarjetavf{Planificación - L4L}{I1}{%
Las empresas usan el sistema Lot-for-Lot (L4L), ya que produce solo lo necesario minimizando los costos totales de inventario.
}{%
\textbf{FALSO.} Porque este sistema asume capacidad ilimitada (lo cual no siempre es verdadero) y asume que los costos de set-up son cero, ignorando el balance con el costo de mantención.
}

\tarjetavf{Gestión de Proyectos}{I1}{%
Las pseudo actividades (o actividades ficticias) son aquellas que requieren más de un recurso a la vez para llevarse a cabo.
}{%
\textbf{FALSO.} Las actividades ficticias no tienen duración ni consumen recursos, solo sirven para mantener la lógica de precedencia en la red.
}

\tarjetavf{Localización}{Mat. Nueva}{%
Tanto el método de análisis de punto de equilibrio como el de centro de gravedad tienen en común que poseen una cantidad acotada de soluciones factibles.
}{%
\textbf{FALSO.} En el método de Centro de Gravedad existe una cantidad infinita de soluciones factibles (cualquier par de coordenadas en el plano continuo).
}

\tarjetavf{Colas y Variabilidad}{I2}{%
El tiempo de flujo promedio es directamente proporcional al porcentaje de utilización de un recurso ($\rho$).
}{%
\textbf{FALSO.} El tiempo de flujo crece de manera \textbf{exponencial} a medida que la utilización ($\rho$) se acerca al 100\%.
}

\tarjetavf{Centro de Distribución}{Mat. Nueva}{%
Las actividades fundamentales que realiza un centro de distribución son: recepcionar, guardar, buscar y despachar.
}{%
\textbf{VERDADERO.} Son los 4 procesos principales (Receiving, Put-away, Picking, Shipping).
}

\tarjetavf{Teoría de Colas}{I2}{%
La ecuación de Kingman o VUT permite determinar de forma analíticamente exacta el tiempo de espera de la cola.
}{%
\textbf{FALSO.} Es una \textbf{aproximación} excelente para sistemas $G/G/1$ u otros distribuidos generalmente, pero no es exacta salvo para distribuciones específicas (ej. $M/M/1$).
}

\tarjetavf{Bodegas - Almacenamiento}{Mat. Nueva}{%
Siempre se debe intentar aumentar las posiciones asignadas en el almacenamiento compartido para mayor eficiencia.
}{%
\textbf{FALSO.} La subdivisión del espacio para compartición está siempre limitada por el aumento en los costos de manipulación, tracking y set-up al separar lotes.
}

\tarjetavf{Calidad}{Mat. Nueva}{%
La calidad percibida por el cliente es netamente objetiva.
}{%
\textbf{FALSO.} Su percepción depende de variables subjetivas: características del producto, las características del servicio, y las \textbf{expectativas y creencias} previas personales del cliente.
}

\tarjetavf{BPMN}{I1}{%
BPMN es mucho mejor que los flujogramas solo porque permite modelar en los procesos actividades complejas y decisiones múltiples.
}{%
\textbf{FALSO.} BPMN es mejor porque además permite modelar explícitamente \textbf{los recursos} y dónde suceden las actividades a través de los \textit{Pools} y \textit{Swimlanes} (carriles).
}

\tarjetavf{Estrategia Operativa}{I1}{%
Un diseño de línea orientado al producto es de flujo continuo y muy eficiente, pero tiene baja flexibilidad.
}{%
\textbf{VERDADERO.} Es altamente automatizado y eficiente para un producto estándar, pero cambiar a otro producto resulta extremadamente costoso y rígido.
}

\tarjetavf{Bodegas - Área Frontal}{Mat. Nueva}{%
Si un SKU es muy popular y se despacha en grandes volúmenes (pallets completos o semi-completos), debe ir en el área de pick frontal.
}{%
\textbf{FALSO.} Un SKU así provocaría un perjuicio al estar en el área frontal, dado que implicaría una pérdida de espacio valioso para otros SKUs que sí requieren pickeos constantes de \textit{menor tamaño}.
}

\tarjetavf{Calidad - Muestreo}{Mat. Nueva}{%
Al disminuir el número o nivel máximo de aceptación en un muestreo (por ej. de 5 unidades a 3 unidades) se aumenta el Riesgo del Productor o error tipo $\alpha$.
}{%
\textbf{VERDADERO.} Al ser más estricto, es más fácil rechazar erróneamente un lote que en realidad sí era bueno (Error Tipo I).
}

\tarjetavf{JIT y Lean}{Mat. Nueva}{%
El objetivo final del Just in Time (JIT) y el Lean es que el sistema productivo logre operar con cero inventario absoluto.
}{%
\textbf{FALSO.} El JIT no busca tener cero inventario, siempre es necesario tener un mínimo inventario transitorio (\textit{WIP} o stock de seguridad pequeño) para proteger contra la variabilidad.
}

\tarjetavf{Six Sigma}{Mat. Nueva}{%
El modelo Six Sigma se enfoca principalmente en la mejora de la calidad de los procesos a través de un estricto control sobre los insumos.
}{%
\textbf{FALSO.} Six Sigma se enfoca principalmente en disminuir la \textbf{variabilidad} del proceso general a través de métodos estadísticos (DMAIC), no solo controlando los insumos de entrada.
}
